\section{Evaluation}
\label{sec: Evaluation}

%Our evaluation aims to answer the following questions to demonstrate the effectiveness of \toolname :
%
%\begin{itemize}[leftmargin=20pt]
%	\item[\textbf{\textit{Q1}}] How many unique and new bugs in real-world DBMSs can \toolname\ detect? (Section~\ref{sec:eval-bugs})
%%	\item[\textbf{\textit{Q2}}] How diverse are the logical bugs found by \toolname? (Section~\ref{sec:eval-diversity})
%	\item[\textbf{\textit{Q2}}] Can \toolname\ find logical bugs that are missed by existing approaches? (Section~\ref{sec:eval-compare})
%\end{itemize}

\subsection{Experimental Setup}
\textbf{Tested DBMS.}
We evaluate \toolname\ on 6 widely used DBMSs, as summarized in Table~\ref{table:testdbms}. 
These DBMSs were chosen based on their popularity and extensive use in real-world applications. 
%MySQL and MariaDB are long-established and widely deployed relational DBMSs, while Percona represents an industrial MySQL-derived system. 
%TiDB is a modern distributed SQL database, and DuckDB and MonetDB are widely recognized analytical DBMSs with strong influence in modern data processing workloads. 
All of these DBMSs have undergone extensive testing for both logical~\cite{hao2023pinolo,kamm2023testing,rigger2020testing,lin2025qtran,rigger2020detecting,rigger2020finding,song2023testing,tang2023detecting,jiang2024detecting,deng2025detecting,zhang2025constant} and crash bugs~\cite{jiang2023dynsql,sqlsmith,pingcap_gorandgen_2022,zhong2020squirrel}.
Although finding new bugs in such widely tested systems is challenging, it serves to validate the effectiveness of \toolname. 
In our experiments, \toolname\ tests the latest versions of each DBMS.  
We run \toolname\ on each DBMS for 24 hours.


\begin{table}[h]
	\centering
	\caption{\textbf{The demographics of the DBMSs under test}}
	\label{table:testdbms}
	\small
	\vspace{-2mm}
	\resizebox{\dimexpr\textwidth/2\relax}{!}{
		\begin{tabular}{lcccc}
			\toprule
			\textbf{DBMS}       & \textbf{Version}           & \textbf{GitHub Stars} & \textbf{DB-Engine}         & \textbf{First Release} \\ 
			\midrule
			MySQL~\cite{mysql}      & 9.5.0             & 11.9K        & 2                 & 1995          \\
			MariaDB~\cite{mariadb}    & 12.2.2 & 6.7K        & 13                & 2009          \\
			MonetDB~\cite{monetdb}  & 11.55.3 & 458 & 150  & 2004          \\
			Percona~\cite{percona}      & 8.0.43-34       & 1.2K         & 122               & 2008          \\
			DuckDB~\cite{duckdb}    & 1.5.0  & 37.3K    & 43               & 2018          \\
			TiDB~\cite{tidb}       & 8.5.5 & 39.4K        & 85                & 2016          \\
			\bottomrule
		\end{tabular}
	}
	\vspace{-2mm}
\end{table}

\textbf{Environment}. 
We conduct the experiments on one server with 104-cores Intel(R) Xeon(R) Gold 6230R CPU @2.10GHz and 500 GB memory. 
The server operates under the Ubuntu 20.04 OS, with the 5.4.0-135-generic version of the Linux kernel.

\subsection{Bug Detection}
\label{sec:eval-bugs}

Table~\ref{tab:bug_summary} summarizes the unique bugs found by \toolname\ across six DBMSs. 
Overall, \toolname\ detected 21 logical bugs, 8 abnormal errors, and 18 crashes, yielding 47 unique bugs in total. 
The logical bugs are distributed across all tested systems, including 5 in MySQL, 3 in MariaDB, 5 in MonetDB, 5 in Percona, 1 in DuckDB, and 2 in TiDB. 
Among all reported bugs, 47 have been confirmed by developers, and 27 have already been fixed. 
These results show that \toolname\ can effectively uncover severe correctness issues even in mature and extensively tested DBMSs.

\begin{table}[t]
	\centering
	\caption{\toolname\ founded 45 unique bugs in six DBMSs.}
	\vspace{-2mm}
	\label{tab:bug_summary}
	\resizebox{\columnwidth}{!}{
		\begin{tabular}{lcccccc}
			\toprule
			\multirow{2}{*}{\textbf{DBMS}} & \multicolumn{3}{c}{\textbf{Bug type}} & \multicolumn{2}{c}{\textbf{Bug status}} \\
			\cmidrule(lr){2-4} \cmidrule(lr){5-6}
			& \textbf{Logical bug} & \textbf{Abnormal error} & \textbf{Crash} & \textbf{Confirmed} & \textbf{Fixed} \\
			\midrule
			MySQL      & 5 & 0 & 3 & 8 & 0 \\
			MariaDB       & 3 & 0 & 0 & 3  & 0 \\
			MonetDB & 5 & 8 & 12 & 25 & 25 \\
			Percona      & 5 & 0 & 1 & 6 & 0 \\
			DuckDB      & 1 & 0 & 2 & 3 & 2 \\
			TiDB        & 2 & 0 & 0 & 2  & 0 \\
			\midrule
			Total       & 21 & 8 & 18 & 47 & 27 \\
			\bottomrule
		\end{tabular}
	}
	\vspace{-3mm}
\end{table}
    
\textbf{Bug Importance.}
Although anecdotal, developer feedback is an important indicator of an approach’s effectiveness and the significance of the bugs found. 
Two DBMS companies reached out to us regarding our testing efforts; both were interested in how we discovered the bugs, and one even invited us to help them find more. 
For example, a developer from PolarDB commented: ``Outstanding work, this tool looks very useful, and we will fully support any technical requirements you may need."
Additionally, we received positive feedback on public bug tracking platforms. 
Among the confirmed bugs, all reported bugs in MariaDB were labeled as ``Major,'' representing the highest severity level, while two bugs in MySQL and two bugs in Percona were classified as ``Serious,'' indicating their significant impact.

\textbf{Throughput.}
We evaluate the efficiency of our test generation by measuring the number of test cases produced per second, where each test case consists of an original DML statement and its corresponding approximate variant. 
Specifically, \toolname\ generates \textcolor{red}{xxx} test cases/sec for MySQL, \textcolor{red}{xxx} for MariaDB, \textcolor{red}{xxx} for MonetDB, \textcolor{red}{xxx} for Percona, \textcolor{red}{xxx} for DuckDB, and \textcolor{red}{xxx} for TiDB. 
Beyond competitive throughput, \toolname\ also attains high execution success rates on most DBMSs, reaching \textcolor{red}{xxx}\% for MySQL, \textcolor{red}{xxx}\% for MariaDB, \textcolor{red}{xxx}\% for MonetDB, \textcolor{red}{xxx}\% for Percona, \textcolor{red}{xxx}\% for DuckDB, and \textcolor{red}{xxx}\% for TiDB.
We further measure the average complexity of generated statements using their size in bytes. 
The average statement size is \textcolor{red}{xxx} bytes for MySQL, \textcolor{red}{xxx} for MariaDB, \textcolor{red}{xxx} for MonetDB, \textcolor{red}{xxx} for Percona, \textcolor{red}{xxx} for DuckDB, and \textcolor{red}{xxx} for TiDB. 
These results show that our structure-guided and constraint-aware generation strategy can efficiently produce a large volume of complex yet executable DML statements.

% Query Generation  （参考 Dinkel）
	% Query Validity (和 DQE 相比)
    % Throughput
    
%\subsection{Comprehensive Bug Analysis}
%\label{sec:eval-diversity}

\textbf{Bug Diversity.}
We analyze the 21 confirmed logical bugs from two perspectives: the diversity of DML types and the structural complexity of the triggering statements. 
In terms of DML types, the majority of bugs are caused by \texttt{UPDATE} statements (16/21), while the remaining involve \texttt{DELETE} (4/21) and \texttt{INSERT} (1/21), indicating that value-modifying operations are particularly error-prone. 
From a structural perspective, most bugs are triggered by complex query constructs: 18 out of 21 bugs involve subqueries, 8 involve joins, 10 involve set operations, and 14 involve logical predicates. 
These results show that DML logical bugs are highly correlated with rich query contexts and interactions among multiple SQL operators, rather than simple statements. 
Such complexity poses a significant challenge to existing approaches like \textsc{DQE}, which are limited to relatively simple and aligned query patterns and cannot effectively support complex constructs such as nested subqueries and set operations. 
In contrast, \toolname\ can systematically generate structurally rich DML statements and leverage state-relation-based oracles to robustly reason about their effects, enabling the detection of these complex and previously overlooked logical bugs.

\begin{table}[t]
	\centering
	\caption{\textbf{Feature distribution of detected DML logical bugs.}}
	\vspace{-2mm}
	\label{tab:bug_pattern}
	\resizebox{\columnwidth}{!}{
		\begin{tabular}{ccccccc}
			\toprule
			\textbf{ID} & \textbf{Bug ID} & \textbf{Type} & \textbf{Subquery} & \textbf{Join} & \textbf{SetOp} & \textbf{LogicOp} \\
			\midrule
			
			1  & \texttt{MySQL\#120047} & UPDATE & \cmark & \cmark & \xmark & \cmark \\
			2  & \texttt{MySQL\#120056} & UPDATE & \cmark & \cmark & \xmark & \cmark \\
			3  & \texttt{MySQL\#120057} & UPDATE & \cmark & \cmark & \cmark & \cmark \\
			4  & \texttt{MySQL\#120058} & UPDATE & \xmark & \cmark & \xmark & \cmark \\
			5  & \texttt{MySQL\#120203} & UPDATE & \cmark & \xmark & \xmark & \xmark \\
			
			6  & \texttt{MariaDB\#39065} & DELETE & \cmark & \xmark & \cmark & \xmark \\
			7  & \texttt{MariaDB\#39097} & UPDATE & \cmark & \xmark & \xmark & \cmark \\
			8  & \texttt{MariaDB\#39144} & UPDATE & \cmark & \xmark & \cmark & \cmark \\
			
			9  & \texttt{MonetDB\#7891} & UPDATE & \cmark & \xmark & \xmark & \cmark \\
			10 & \texttt{MonetDB\#7888} & UPDATE & \xmark & \xmark & \xmark & \cmark \\
			11 & \texttt{MonetDB\#7889} & DELETE & \cmark & \xmark & \cmark & \xmark \\
			12 & \texttt{MonetDB\#7890} & INSERT & \cmark & \xmark & \cmark & \cmark \\
			13 & \texttt{MonetDB\#7897} & DELETE & \cmark & \xmark & \cmark & \cmark \\
			
			14 & \texttt{Percona\#10877} & UPDATE & \cmark & \cmark & \xmark & \cmark \\
			15 & \texttt{Percona\#10878} & UPDATE & \cmark & \cmark & \xmark & \cmark \\
			16 & \texttt{Percona\#10879} & UPDATE & \cmark & \cmark & \cmark & \cmark \\
			17 & \texttt{Percona\#10906} & UPDATE & \xmark & \cmark & \xmark & \cmark \\
			18 & \texttt{Percona\#11008} & UPDATE & \cmark & \xmark & \xmark & \xmark \\
			
			19 & \texttt{DuckDB\#21609} & INSERT & \cmark & \cmark & \xmark & \cmark \\
			
			20 & \texttt{TiDB\#67018} & UPDATE & \cmark & \xmark & \cmark & \xmark \\
			21 & \texttt{TiDB\#67019} & DELETE & \cmark & \xmark & \cmark & \cmark \\
			
			\bottomrule
		\end{tabular}
	}
	\vspace{-3mm}
\end{table}

\textbf{Case Study.}

% RQ3：Comoporehensive Bug Analysis
	% Bug Statements：Update、Delete、Insert (Rules) 一个大表  id，bug id，type，（rule），EET的feature
% Selected Bugs
%   Update 、 Delete 、 Insert 各一个

\subsection{Comparative Study}
\label{sec:eval-compare}
To demonstrate that \toolname\ can detect DML logical bugs missed by existing methods, we conduct our comparative study from two perspectives: tool level and oracle level. 
At the tool level, we compare \toolname\ with six representative logical bug detection approaches, including \textsc{PQS}~\cite{rigger2020testing}, \textsc{NoREC}~\cite{rigger2020detecting}, \textsc{TLP}~\cite{rigger2020finding}, \textsc{DQE}~\cite{song2023testing}, \textsc{Pinolo}~\cite{hao2023pinolo}, and \textsc{EET}~\cite{jiang2024detecting}, by running each tool on the latest version of each supported DBMS for 24 hours. 
Since the baseline methods do not support all DBMSs tested by \toolname, the comparison is restricted to the DBMSs commonly supported by \toolname\ and each baseline method. 
As shown in Table~\ref{tab:tool-support-dbms}, \toolname\ consistently outperforms the baselines in terms of triggered logical bugs, finding 5, 3, 1, 10, 7, and 4 more bugs than \textsc{PQS}, \textsc{NoREC}, \textsc{TLP}, \textsc{DQE}, \textsc{Pinolo}, and \textsc{EET}, respectively, where each increment is computed only on the commonly supported DBMSs. Overall, \toolname\ triggers 21 logical bugs within 24 hours, whereas all baseline tools together trigger only 6.
At the oracle level, we analyze whether existing testing oracles are capable of detecting the logical bugs uncovered by \toolname. 
From a theoretical perspective, the \logical\ logical bugs detected by \toolname\ fundamentally rely on reasoning about state transitions, including which rows are affected and how their values are modified. 
However, existing query-centric approaches, such as \textsc{PQS}, \textsc{NoREC}, \textsc{TLP}, \textsc{Pinolo}, and \textsc{EET}, define correctness over returned query results and therefore cannot capture errors that only manifest in persistent state changes. 
Although \textsc{DQE} targets DML statements, its oracle is based on cross-statement row-access consistency and mainly checks whether related statements access the same tuples, making it insufficient for detecting value-level errors and many complex state transitions induced by joins, subqueries, and other rich DML constructs. 
As a result, the logical bugs detected by \toolname\ are not observable under the oracles used by these baseline approaches. 
This explains why the \logical\ logical bugs detected by \toolname\ are entirely disjoint from the 6 logical bugs triggered by the baselines, and highlights that \toolname\ explores a different and previously untested space of DML correctness. 
%At the oracle level, we further analyze why \toolname\ can uncover bugs missed by existing tools.  
%The \logical\ logical bugs detected by \toolname\ are entirely disjoint from the 6 logical bugs triggered by the baselines, indicating that \toolname\ discovers a different class of DML bugs. 
%This advantage stems from its stronger oracle design. 
%Existing query-centric approaches, such as \textsc{PQS}, \textsc{NoREC}, \textsc{TLP}, \textsc{Pinolo}, and \textsc{EET}, define correctness over returned query results and therefore are not naturally suited for reasoning about persistent state changes caused by DML execution, while \textsc{DQE} relies on cross-statement row-access consistency and mainly checks whether related statements touch the same tuples. 
%As a result, these approaches cannot naturally validate many value-level errors and complex state transitions induced by joins, subqueries, and other rich DML constructs. 
%In contrast, \toolname\ reasons directly about DML semantics through state effects and checks approximate relations over state changes, enabling it to detect logical bugs that are invisible to result-set-based or row-access-based oracles.
%These results show that \toolname\ not only outperforms existing methods in practice, but also complements them by exposing previously overlooked DML logical bugs.

%To compare \toolname\ against state-of-the-art DBMS testing methods, we selected six representative logical bug detection approaches for comparison: \textsc{PQS}~\cite{rigger2020testing}, \textsc{NoREC}~\cite{rigger2020detecting}, \textsc{TLP}~\cite{rigger2020finding}, \textsc{DQE}~\cite{song2023testing}, \textsc{Pinolo}~\cite{hao2023pinolo}, and \textsc{EET}~\cite{jiang2024detecting}. 
%We ran each tool on the latest version of each supported DBMS for 24 hours. 
%Since the baseline methods do not support all DBMSs tested by \toolname, the comparison is restricted to the DBMSs commonly supported by \toolname\ and each baseline method. 
%As shown in Table~\ref{tab:tool-support-dbms}, \toolname\ consistently outperforms the other tools in detecting logical bugs across the evaluated DBMSs. 
%Specifically, \toolname\ found 5, 3, 1, 10, 7, and 4 more bugs than \textsc{PQS}, \textsc{NoREC}, \textsc{TLP}, \textsc{DQE}, \textsc{Pinolo}, and \textsc{EET}, respectively. 
%Note that these increments are computed only on the DBMSs commonly supported by both approaches. 
%We further analyzed the relation between the bugs detected by existing tools and those detected by \toolname\ to examine whether \toolname\ can uncover bugs missed by prior methods. 
%Our results show that the \reported\ logical bugs detected by \toolname\ are entirely disjoint from the 6 logical bugs triggered by the baseline tools, indicating that \toolname\ can discover previously missed bugs and effectively complements existing methods.

\begin{table}[t]
	\centering
	\vspace{-2mm}
	\caption{\textbf{Number of logical bugs triggered in 24 hours.}}
	\vspace{-2mm}
	\small
	\label{tab:tool-support-dbms}
	\resizebox{\columnwidth}{!}{
		\begin{tabular}{l|ccccccc}
			\toprule
			\textbf{DBMS} & \textbf{PQS} & \textbf{NoREC} & \textbf{TLP} & \textbf{DQE} & \textbf{Pinolo} & \textbf{EET} & \textbf{\toolname} \\
			\midrule
			MySQL    & 0 & - & 4 & 0 & 2 & 3 & 5 \\
			MariaDB  & - & 0 & - & 0 & 1 & - & 3 \\
			MonetDB  & - & - & - & - & - & - & 5 \\
			Percona  & - & - & - & - & - & - & 5 \\
			DuckDB   & - & - & - & - & - & - & 1 \\
			TiDB     & - & - & 2 & 0 & 0 & 0 & 2 \\
			\midrule
			\textbf{Total}     & 0 & 0 & 6 & 0 & 3 & 3 & 21 \\
			\textbf{Increment} & 5 & 3 & 1 & 10 & 7 & 4 &  \\
			\bottomrule
		\end{tabular}
	}
	\raggedright
	\noindent \textbf{Note.} ``-'' indicates that the tool does not support the DBMS. \\
	``Increment'' is computed on the commonly supported DBMSs between the baseline tool and \toolname.
	\vspace{-5mm}
\end{table}

%\textbf{Temporal Comparison.}
%We also present a temporal comparison to evaluate whether \toolname\ can uncover DML logical bugs that have long existed but were missed by prior approaches. 
%For each confirmed logical bug discovered by \toolname, we trace its earliest bug-involved version and compare it against the publication timeline of existing DBMS testing approaches~\cite{rigger2020detecting,rigger2020finding,rigger2020testing,song2023testing,tang2023detecting,jiang2024detecting,song2024detecting,hao2023pinolo,deng2025detecting,ba2024keep}. 
%If a bug had already existed before the release of prior techniques but remained undiscovered until \toolname, we regard it as evidence that previous approaches failed to detect it. 
%This comparison is reasonable for two reasons. 
%First, the bugs reported by \toolname\ were not marked as duplicates by the corresponding developers, indicating that they had not been previously reported. 
%Second, most of the evaluated DBMSs have already been extensively tested by existing approaches, yet these long-lived DML bugs still remained undiscovered. 
%
%As shown in Table~\ref{tab:bug-latency}, among the confirmed logical bugs found by \toolname, \textcolor{red}{xxx} bugs can be traced back to versions released before 2023, and \textcolor{red}{xxx} of them were already present before 2020. 
%This means that a substantial number of DML logical bugs had survived multiple generations of prior DBMS testing techniques. 
%The main reason is that previous approaches focus either on query-result relations or on cross-statement row-access consistency, and therefore cannot effectively reason about the correctness of complex DML-induced state transitions. 
%In contrast, \toolname\ directly targets DML semantics by combining complex DML generation with state-relation-based metamorphic oracles, making it particularly effective at exposing long-standing logical bugs in state-modifying statements. 
%We further observe that \toolname\ identifies at least one long-lived bug in each tested DBMS, with the longest latency reaching \textcolor{red}{xxx} years. 
%These results demonstrate that \toolname\ is capable of uncovering long-standing, previously overlooked DML logical bugs across diverse DBMSs.
%
%\begin{table}[t]
%	\centering
%	\caption{\textbf{Latency of the confirmed logical bugs.}}
%	\vspace{-2mm}
%	\label{tab:bug-latency}
%	\small
%	\begin{tabular}{lcccc}
%		\toprule
%		\multirow{2}{*}{\textbf{DBMS}} &
%		\multirow{2}{*}{\textbf{Found}} &
%		\multicolumn{2}{c}{\textbf{Bug-involved year}} &
%		\multirow{2}{*}{\textbf{Longest latency}} \\
%		\cmidrule(lr){3-4}
%		& & \textbf{< 2023} & \textbf{< 2020} & \\
%		\midrule
%		MySQL    & \textcolor{red}{xxx} & \textcolor{red}{xxx} & \textcolor{red}{xxx} & \textcolor{red}{xxx} years \\
%		MariaDB  & \textcolor{red}{xxx} & \textcolor{red}{xxx} & \textcolor{red}{xxx} & \textcolor{red}{xxx} years \\
%		MonetDB  & \textcolor{red}{xxx} & \textcolor{red}{xxx} & \textcolor{red}{xxx} & \textcolor{red}{xxx} years \\
%		Percona  & \textcolor{red}{xxx} & \textcolor{red}{xxx} & \textcolor{red}{xxx} & \textcolor{red}{xxx} years \\
%		DuckDB   & \textcolor{red}{xxx} & \textcolor{red}{xxx} & \textcolor{red}{xxx} & \textcolor{red}{xxx} years \\
%		TiDB     & \textcolor{red}{xxx} & \textcolor{red}{xxx} & \textcolor{red}{xxx} & \textcolor{red}{xxx} years \\
%		\midrule
%		\textbf{Total} & \textcolor{red}{xxx} & \textcolor{red}{xxx} & \textcolor{red}{xxx} & \textcolor{red}{xxx} years \\
%		\bottomrule
%	\end{tabular}
%	\vspace{-2mm}
%\end{table}

% RQ4： Comparative Study
% Branch Coverage、Code Coverage (Optional)
% 和 SOSP 的一样吧， 加一个 Venn图（参考 EDC）



